\section{Le problème et l'objectif}\label{sec-probleme}

\begin{objectif}
Décrire la chaîne qui va des contributions Totem aux lignes de contrôle, montrer pourquoi une courbe de repo implicite
paraît souvent aberrante sans l'être, et préciser ce que la note doit livrer.
\end{objectif}

\subsection{La chaîne de contrôle}
Pour chaque sous-jacent et chaque date de contrôle, le contrôleur dispose de trois sources (\cref{fig-chaine}).
\begin{itemize}
\item \textbf{Totem options} : pour chaque maturité $T_k$, des prix de consensus de calls et de puts sur une grille
de strikes, le nombre de contributeurs $n^{\opt}_k$ et la dispersion de leurs cotations.
\item \textbf{Totem dividendes} : pour chaque période $p$ (en général une année civile), le niveau de consensus
$\ell^{\tot}_p$, c'est-à-dire la somme des dividendes détachés dans la période, le nombre de contributeurs
$n^{\dv}_p$ et leur dispersion.
\item \textbf{Le front office} : le spot, la courbe des taux, la courbe de repo $q^{\fo}$, la prévision de dividendes
(dates de détachement et montants) et les sensibilités de chaque portefeuille au repo et aux dividendes.
\end{itemize}
La parité call--put donne les forwards Totem $F^{\tot}_k$. Les niveaux de dividendes sont répartis sur les dates de
détachement du FO. Un bootstrap donne ensuite la courbe de repo $q^{\tot}$ qui, avec ces dividendes, reproduit
exactement les forwards Totem. Cette courbe sert à deux choses : calculer les lignes de contrôle, c'est-à-dire les
écarts Totem moins FO multipliés par les sensibilités (\cref{sec-controle}), et contrôler la qualité de la courbe
elle-même (\cref{sec-detection}).

\begin{figure}[ht]
\centering
\begin{tikzpicture}[
  font=\small, >={Stealth[length=2mm]},
  box/.style={rounded corners=2pt, align=center, inner sep=4pt, minimum height=1.2cm},
  src/.style={box, draw=accent, fill=accentwash, minimum width=3.3cm},
  fo/.style={box, draw=muted, fill=greywash, minimum width=3.3cm},
  mid/.style={box, draw=rule, fill=white, minimum width=2.2cm},
  key/.style={box, draw=accentdark, line width=0.9pt, fill=white, minimum width=2.2cm},
  use/.style={box, draw=green, fill=greenwash, minimum width=3.5cm},
  lab/.style={font=\footnotesize\itshape, text=ink2, align=center}]
\node[src] (opt) at (0,2.4) {\textbf{Totem options}\\ calls et puts, $n^{\opt}_k$};
\node[src] (div) at (0,0.8) {\textbf{Totem dividendes}\\ niveaux $\ell^{\tot}_p$, $n^{\dv}_p$};
\node[fo] (fo) at (0,-1.0) {\textbf{Front office}\\ $q^{\fo}$, dividendes,\\ sensibilités $s^q$, $s^\ell$};
\node[mid] (F) at (4.6,2.4) {forwards\\ $F^{\tot}_k$};
\node[mid] (D) at (4.6,0.8) {dividendes\\ $D^{\tot}_k$};
\node[key] (Q) at (7.65,1.6) {\textbf{repo implicite}\\ $q^{\tot}_k$};
\node[use] (L) at (11.25,2.4) {\textbf{Lignes de contrôle}\\ écarts $\times$ sensibilités\\ \cref{sec-controle}};
\node[use] (Z) at (11.25,0.6) {\textbf{Détection}\\ sur l'écart de portage $\delta_k$\\ \cref{sec-detection}};
\draw[->] (opt) -- node[lab, above] {parité} (F);
\draw[->] (div) -- node[lab, above] {allocation} (D);
\draw[->] (F.east) -- (Q);
\draw[->] (D.east) -- (Q);
\node[lab, anchor=south] at ($(Q.north)+(0,0.05)$) {bootstrap};
\draw[->] (Q) -- (L);
\draw[->] (Q) -- (Z);
\draw[->, dashed, draw=muted] ($(fo.north east)+(0,-0.15)$) -- (D.south);
\node[lab, anchor=north west] at ($(D.south)+(0.1,-0.12)$) {dates de détachement};
\draw[draw=muted] (fo.east) -- (13.45,-1.0) -- (13.45,2.4);
\draw[->, draw=muted] (13.45,2.4) -- (L.east);
\draw[->, draw=muted] (13.45,0.6) -- (Z.east);
\node[lab, anchor=north] at (7.65,-1.05) {$q^{\fo}$, dividendes FO, sensibilités};
\end{tikzpicture}
\caption{La chaîne de contrôle. Les deux services Totem ont des contributeurs différents, et leur nombre varie avec la
maturité ; la courbe de repo implicite $q^{\tot}$ hérite du bruit des deux.}
\label{fig-chaine}
\end{figure}

\subsection{Pourquoi la courbe paraît aberrante}
Deux raisons, qui n'ont rien à voir avec une erreur, suffisent à donner à une courbe saine un aspect irrégulier.

\paragraph{Deux panels, deux moyennes.} Le forward Totem est la moyenne des forwards du panel options, et le forward
de chaque contributeur contient sa propre vue des dividendes. Le contrôleur lui retire les dividendes d'un autre panel,
celui du service dividendes. Si les deux panels n'ont pas la même vue moyenne des dividendes, la différence se
retrouve dans le repo implicite. Ce n'est pas une erreur de Totem : c'est la conséquence mécanique de deux moyennes
prises sur des contributeurs différents. Sa taille dépend des nombres de contributeurs, qui baissent avec la maturité,
et de la dispersion des vues sur les dividendes, qui augmente avec elle (\cref{sec-gap}).

\paragraph{L'effet loupe en $1/T$.} Le repo est un taux : une erreur $\varepsilon$ sur le logarithme du forward devient
une erreur $\varepsilon/T$ sur le repo. Une erreur de \pct{0.1} sur le forward vaut \num{120}\,pb de repo à un mois et
\num{2}\,pb à cinq ans. Le court terme paraît donc toujours bruité, tandis qu'une erreur importante sur le long terme
passe inaperçue.

Une courbe d'aspect irrégulier n'est donc pas forcément mauvaise, et une courbe d'aspect lisse n'est pas forcément
bonne. Toute la suite travaille sur l'\emph{écart de portage} $\delta_k=(q^{\tot}_k-q^{\fo}_k)\,T_k$, où les erreurs
s'additionnent sans être amplifiées, et le compare à une bande de bruit calculée à partir des statistiques que Totem
publie.

\subsection{L'exemple fil rouge}\label{sec-filrouge}
\begin{filrouge}
Une action européenne synthétique (\cref{sec-banc}) : spot $S_0=100$, un dividende annuel d'environ \vRunDiv{}
détaché début mai (le premier le \vRunExDate), un repo FO qui passe de \vRunQfoShort\,pb à court terme à
\vRunQfoLong\,pb à cinq ans. Totem compte \vRunNoptShort{} contributeurs options sur les maturités courtes et
\vRunNoptLong{} à cinq ans ; le service dividendes en compte \vRunNdivShort{} sur 2027 et \vRunNdivLong{} sur 2031.
Les onze maturités Totem vont d'un mois à cinq ans. La \cref{fig-clean} montre le marché sain ; le même sous-jacent
est repris, avec le même bruit de contributeurs, sous huit anomalies en \cref{sec-signatures} (\cref{fig-gallery}) et
passé au détecteur en \cref{sec-retour}.
\end{filrouge}

\begin{figure}[ht]
\centering\includegraphics[width=\linewidth]{figures/\lang/clean.pdf}
\caption{Exemple fil rouge, marché sain. (a) Repo FO et repo implicite Totem, avec la bande de $\pm2$ écarts-types du
bruit de consensus autour du repo FO. (b) Le même écart vu en portage : chaque point est
$\delta_k=(q^{\tot}_k-q^{\fo}_k)\,T_k$, en points de base du spot.}
\label{fig-clean}
\end{figure}

Dans l'espace de repo, la courbe Totem paraît mal formée : elle perd \vRunDrop\,pb entre six et neuf mois et finit
\vRunGapLong\,pb sous le FO à cinq ans. Dans l'espace de portage, tous les points restent dans la bande.
Le décrochage à neuf mois correspond à un écart de \vRunDivGap{} sur le dividende de mai 2027 entre la vue du panel
options et celle du panel dividendes, soit \vRunDivGapPct{} du dividende et \vRunZNine{} écart-type ; à cinq ans,
l'écart de portage vaut \vRunDeltaFive\,pb du spot pour un écart-type de \vRunSdFive\,pb. Le détecteur de la
\cref{sec-detection} ne signale rien. Cette courbe est saine : c'est l'aspect normal d'un repo implicite tiré de deux
panels différents.

\subsection{Ce que la note doit livrer}
Les trois objectifs de la première page se traduisent en quatre livrables.
\begin{enumerate}
\item \textbf{Un diagnostic} (O1, \cref{sec-diagnostic}) : une décomposition du repo implicite qui sépare le bruit
attendu, dont on sait calculer la taille, des mécanismes qui produisent de vraies anomalies, chacun avec une forme
reconnaissable.
\item \textbf{Une règle de contrôle} (O2, \cref{sec-controle}) : une décomposition des lignes de contrôle qui ne dépend
pas des contributeurs dividendes là où c'est possible, et une mesure, calculée à partir des sensibilités, de ce
qu'elle laisse dépendre d'eux.
\item \textbf{Un détecteur} (O3, \cref{sec-detection}) dont le taux de fausses alertes est connu, qui nomme le mécanisme
en cause, en donne l'amplitude en unités naturelles (spot décalé de tant, dividende de mai 2027 trop bas de tant),
répare la courbe et distingue une anomalie Totem d'une erreur de marquage FO.
\item \textbf{Une comparaison} (\cref{sec-banc}) avec les idées de départ : coefficients B-spline par région et
catégorie regroupés par mélange gaussien, petits réseaux de neurones spécialisés par type d'anomalie.
\end{enumerate}

\begin{retenir}
\begin{itemize}[leftmargin=12pt]
\item La courbe de repo implicite est un résidu : forwards d'un panel moins dividendes d'un autre panel. Elle hérite du
bruit des deux.
\item L'œil trompe : le repo amplifie en $1/T$ les erreurs du court terme et écrase celles du long terme. On juge la
courbe sur l'écart de portage $\delta_k$, avec une bande de bruit.
\item Dans l'exemple fil rouge, une courbe qui perd \vRunDrop\,pb en trois mois et finit \vRunGapLong\,pb sous
le FO est parfaitement saine.
\end{itemize}
\end{retenir}
